% ============================================================
% 5.0 Tổng quan hiện thực
% ============================================================

\section{Tổng quan hiện thực}
\label{sec:version-overview}

Phiên bản hiện tại của ShopNexus (\textbf{v1.0 Development}) đã hiện thực toàn bộ luồng nghiệp vụ cốt lõi của nền tảng thương mại điện tử, bao gồm: quản lý tài khoản và xác thực JWT, catalog sản phẩm SPU/SKU với tìm kiếm lai (Milvus + PostgreSQL), giỏ hàng và checkout với reserve inventory, thanh toán qua VNPay và COD, quản lý đơn hàng cho buyer/seller, hoàn trả và tranh chấp, khuyến mãi tự động, chat real-time, phân tích hành vi người dùng, và dashboard thống kê cho seller.

\subsection{Công nghệ sử dụng}
\label{subsec:tech-stack}

Bảng công nghệ chi tiết đã trình bày tại Chương~2. Tóm tắt các thành phần chính:

\begin{table}[htbp]
\centering
\caption{Tóm tắt công nghệ chính}
\label{tab:tech-summary}
\begin{tabular}{@{}llp{7cm}@{}}
\toprule
\rowcolor{headerblue}
\thead{Tầng} & \thead{Công nghệ} & \thead{Vai trò} \\
\midrule
Backend & Go 1.26, Echo v4, Restate & HTTP server, durable execution \\
Database & PostgreSQL 18, Redis 8, Milvus & RDBMS, cache, vector search \\
Storage & MinIO, NATS JetStream & Object storage, message queue \\
External & VNPay, GHTK, Nominatim & Thanh toán, vận chuyển, geocoding \\
Frontend & Next.js 16, React 19, shadcn/ui & Giao diện người dùng \\
\bottomrule
\end{tabular}
\end{table}

\subsection{Triển khai}
\label{subsec:deployment}

Môi trường phát triển sử dụng \textbf{Docker Compose} để quản lý hạ tầng gồm: PostgreSQL, Redis, cụm Restate 3 node, Milvus (tìm kiếm vector), và MinIO (lưu trữ đối tượng). Server Go chạy trên máy host với hot-reload, kết nối các service qua port mapping. Quy trình khởi tạo: khởi động container $\rightarrow$ chạy migration $\rightarrow$ nạp dữ liệu mẫu $\rightarrow$ khởi động server $\rightarrow$ đăng ký service với Restate.

% ============================================================
\section{Chiến lược kiểm thử}
\label{sec:test-strategy}

Phần này trình bày chiến lược kiểm thử toàn diện cho ShopNexus theo mô hình \textbf{Test Pyramid} --- nguyên tắc SQA nhằm tối ưu tỷ lệ giữa chi phí test và độ tin cậy mà nó cung cấp.

\subsection{Mô hình Test Pyramid}
\label{subsec:test-pyramid}

\begin{center}
\begin{tikzpicture}[
  layer/.style={trapezium, trapezium left angle=70, trapezium right angle=110, draw=headerblue, align=center, font=\footnotesize},
  lbl/.style={font=\scriptsize, text=gray!50!black, align=left, anchor=west},
]
\node[layer, fill=red!30,    minimum width=3cm, minimum height=0.8cm] (e2e)  at (0, 3.5) {E2E / Chaos\\\scriptsize 10\%};
\node[layer, fill=orange!40, minimum width=5.5cm, minimum height=0.9cm] (int)  at (0, 2.3) {Integration\\\scriptsize 20\%};
\node[layer, fill=headerblue!30, minimum width=8.5cm, minimum height=1.2cm] (unit) at (0, 0.8) {Unit tests\\\scriptsize 70\%};

\node[lbl] at (4.7, 3.5) {Ít --- chậm --- đắt};
\node[lbl] at (4.7, 2.3) {Trung bình};
\node[lbl] at (4.7, 0.8) {Nhiều --- nhanh --- rẻ};
\end{tikzpicture}
\captionof{figure}{Test Pyramid --- phân bổ tỷ lệ test theo cấp độ}
\label{fig:test-pyramid}
\end{center}

\textbf{Tỷ lệ mục tiêu}: \textbf{70\% unit / 20\% integration / 10\% E2E}. Unit test nhanh (mili-giây), chạy mỗi lần commit qua \texttt{go test}; integration test vừa phải (giây), chạy trước khi merge; E2E test chậm (chục giây đến phút), chạy nightly và pre-release.

\subsection{Phân loại test theo cấp độ}
\label{subsec:test-types}

\begin{table}[H]
\centering
\caption{Phân loại test áp dụng cho ShopNexus}
\label{tab:test-types}
\small
\begin{tabular}{@{}p{2.8cm}p{4.5cm}p{3cm}p{3cm}@{}}
\toprule
\rowcolor{headerblue}
\thead{Loại test} & \thead{Phạm vi kiểm thử} & \thead{Công cụ} & \thead{Ví dụ} \\
\midrule
Unit test & Một hàm, không truy cập DB / mạng & Go test + table-driven & Test serialization, test chuyển trạng thái vận chuyển \\
Integration test & Tầng nghiệp vụ + DB thật (container), không UI & dockertest + PostgreSQL & Luồng checkout đầy đủ: đặt chỗ $\to$ tạo mục $\to$ xác nhận $\to$ tạo đơn \\
Contract test & Giao tiếp liên module qua proxy & Restate test harness & Tra cứu SKU từ Catalog trả đúng cấu trúc \\
E2E test & Toàn bộ: HTTP $\to$ nghiệp vụ $\to$ DB $\to$ nhà cung cấp & Postman / k6 & Thanh toán VNPay sandbox thành công \\
Load / Performance & Hệ thống chịu tải N người dùng đồng thời & k6, vegeta & 100 checkout đồng thời, đo p99 latency \\
Chaos / Fault Injection & Hệ thống phục hồi khi thành phần bị lỗi & pumba kill container & Dừng node giữa luồng checkout $\to$ replay tự động \\
Security / Pen-test & Khai thác lỗ hổng chủ động & gosec, OWASP ZAP, thủ công & SQL injection, JWT bypass, giả mạo webhook \\
Static Analysis & Lỗi biên dịch + chất lượng mã nguồn & golangci-lint, staticcheck & Phụ thuộc vòng, gán không hiệu quả, mã không dùng \\
Mutation test & Đo độ mạnh của bộ test & go-mutesting & Đổi \texttt{<} thành \texttt{>} trong công thức $\to$ test phải fail \\
\bottomrule
\end{tabular}
\end{table}

\subsection{Tiêu chí chất lượng theo mức phát triển}
\label{subsec:quality-levels}

Mỗi chức năng cốt lõi của ShopNexus được kiểm soát chất lượng qua ba mức: \textit{ý niệm} (phân tích --- ``feature cần gì?''), \textit{thiết kế} (kiến trúc --- ``code xử lý như thế nào?''), và \textit{hiện thực} (kiểm thử --- ``chứng minh bằng test nào?''). Bảng~\ref{tab:quality-levels} trình bày đầy đủ cho 7 chức năng, bao gồm cả các tính năng bổ sung gần đây (dispute, transport webhook, partial refund).

\rowcolors{2}{white}{rowgray}
\begin{longtable}{@{}p{2.8cm}p{3.8cm}p{3.8cm}p{3.8cm}@{}}
\caption{Tiêu chí chất lượng theo ba mức phát triển}\label{tab:quality-levels}\\
\toprule
\rowcolor{headerblue}
\thead{Chức năng} & \thead{Ý niệm (Phân tích)} & \thead{Thiết kế} & \thead{Hiện thực (Test)} \\
\midrule
\endfirsthead
\toprule
\rowcolor{headerblue}
\thead{Chức năng} & \thead{Ý niệm (Phân tích)} & \thead{Thiết kế} & \thead{Hiện thực (Test)} \\
\midrule
\endhead
\midrule
\multicolumn{4}{r}{\small\textit{Tiếp theo trang sau}} \\
\endfoot
\bottomrule
\endlastfoot
\textbf{Đặt hàng} &
Snapshot giá tại thời điểm đặt; đặt chỗ kho nguyên tử; xử lý exactly-once &
API idempotent; gọi liên module qua proxy; mọi side effect trong giao dịch durable &
Integration test replay journal; test tồn kho sau crash; test serialization schema \\

\textbf{Thanh toán} &
Bảo mật; exactly-once (không trừ tiền 2 lần dù webhook gửi lại) &
Xác minh chữ ký HMAC webhook; thời hạn phiên thanh toán &
Test webhook hợp lệ / giả mạo; test timeout phiên; test replay callback \\

\textbf{Xác nhận đơn} &
Phân quyền: chỉ seller sở hữu sản phẩm; tính tổng chính xác &
Kiểm tra quyền sở hữu tại tầng nghiệp vụ; công thức tổng tiền thống nhất &
Negative test: seller xác nhận đơn seller khác (phải từ chối); test công thức \\

\textbf{Hoàn trả} &
Chỉ đơn đã thanh toán thành công; tối đa 1 hoàn trả đang xử lý/đơn; hỗ trợ hoàn trả một phần &
Kiểm tra hoàn trả trùng trước khi tạo; liên kết bằng chứng qua module tài nguyên &
Test hoàn trả trùng bị từ chối; test tạo 2 hoàn trả đồng thời cùng đơn \\

\textbf{Tranh chấp} &
Buyer hoặc seller mở tranh chấp; hoàn trả chưa kết thúc mới được tranh chấp &
Kiểm tra người gọi là buyer/seller của đơn; kiểm tra không có tranh chấp đang chờ; giới hạn 3 lần/phút &
Test phân quyền: người ngoài mở tranh chấp $\to$ 403; test giới hạn tốc độ $\to$ 429 \\

\textbf{Giao vận} &
Trạng thái giao vận phản ánh đúng thực tế từ webhook nhà vận chuyển &
State machine: chỉ cho phép chuyển trạng thái hợp lệ; trạng thái cuối không chuyển tiếp &
16 subtests chuyển trạng thái; test bỏ qua trạng thái (Pending $\to$ InTransit) bị từ chối \\

\textbf{Dashboard} &
Phản hồi dưới 1 giây cho trải nghiệm người dùng &
Partial index + SQL aggregate, không xử lý vòng lặp ở ứng dụng &
Benchmark 10K+ đơn hàng; đo thời gian phản hồi dưới tải \\
\end{longtable}
\rowcolors{1}{}{}

\subsection{Mục tiêu coverage và chất lượng}
\label{subsec:quality-targets}

\begin{table}[htbp]
\centering
\caption{Mục tiêu chất lượng code và test}
\label{tab:quality-targets}
\small
\begin{tabular}{@{}p{4cm}p{3.5cm}p{6cm}@{}}
\toprule
\rowcolor{headerblue}
\thead{Chỉ số} & \thead{Mục tiêu} & \thead{Cách đo} \\
\midrule
Line coverage (unit test) & $\geq$ 70\% & \texttt{go test -cover ./...} \\
Branch coverage (critical path) & $\geq$ 60\% & \texttt{go test -covermode=count} \\
Cyclomatic complexity & $\leq$ 10 / function & \texttt{gocyclo} \\
LCOM (lack of cohesion) & $\leq$ 0.5 & Phân tích tĩnh \\
Function length & $\leq$ 80 dòng & Review guideline \\
File length & $\leq$ 500 dòng & Soft limit (split nếu vượt) \\
Mutation score (critical) & $\geq$ 60\% & \texttt{go-mutesting} trên module order/biz \\
Số lỗi nghiêm trọng (\texttt{gosec}) & 0 & CI bắt buộc pass \\
\bottomrule
\end{tabular}
\end{table}

\subsection{Tiêu chí kiểm thử hiện thực}
\label{subsec:impl-test-criteria}

\begin{table}[htbp]
\centering
\caption{Tiêu chí kiểm thử cho phần hiện thực}
\label{tab:impl-test-criteria}
\small
\begin{tabular}{@{}p{2.5cm}p{5cm}p{3.5cm}p{2.8cm}@{}}
\toprule
\rowcolor{headerblue}
\thead{Tiêu chí} & \thead{Phép thử / Test-inputs} & \thead{Kết quả mong đợi} & \thead{Cấp độ} \\
\midrule
Xử lý chính xác & Black box: test input hợp lệ / không hợp lệ; White box: phủ branch của biz logic & Đạt spec thiết kế; tất cả branch quan trọng được test & Unit + Integration \\
Hiệu năng & Stress test: 100 concurrent checkout, 500 rps read; đo p50/p95/p99 & p99 read $\leq$ 500ms, p99 write $\leq$ 2s; không có memory leak & Load test \\
Bảo mật & Code review (\texttt{gosec}); pen-test: SQL injection, JWT forge, webhook replay & Không lỗi nghiêm trọng; mọi endpoint reject tấn công & Security test \\
Phục hồi & Chaos: kill Restate node hoặc DB giữa durable step; quan sát replay & State nhất quán sau replay; không double-side-effect & Chaos test \\
Regression & Full test suite pre-release; so sánh coverage với baseline & Không giảm coverage; 0 test mới fail & CI gate \\
\bottomrule
\end{tabular}
\end{table}

\subsection{Quy trình kiểm thử trong CI/CD}
\label{subsec:ci-testing}

Mọi Pull Request đều chạy qua pipeline sau trước khi được merge:

\begin{enumerate}
    \item \textbf{Static analysis}: \texttt{golangci-lint run} + \texttt{staticcheck ./...} + \texttt{gosec ./...} --- dừng nếu có lỗi nghiêm trọng.
    \item \textbf{Unit tests}: \texttt{go test -race -cover ./...} --- race detector bật để bắt concurrency bug; coverage không được giảm.
    \item \textbf{Integration tests}: spawn PostgreSQL / Redis / Restate qua \texttt{dockertest}; chạy test hit DB thật.
    \item \textbf{Build}: \texttt{go build ./...} --- đảm bảo compile sạch trên mọi nhánh code.
    \item \textbf{Peer review}: ít nhất 1 reviewer approve; áp dụng checklist thiết kế (SOLID, testability, modularity) cho file mới.
\end{enumerate}

Các loại test nặng (E2E, load, chaos, mutation) chạy theo lịch: E2E nightly trên staging, load test mỗi release, mutation test mỗi sprint cho module thay đổi nhiều.
